\documentclass[11pt]{article}
\newcommand{\ListaTitulo}{Gabarito 09: VA discreta, esperança, variância e FDA}
% Preâmbulo compartilhado: MATF14, Listas de Exercícios 2026
% Usado por ListaNN.tex e GabaritoNN.tex via % Preâmbulo compartilhado: MATF14, Listas de Exercícios 2026
% Usado por ListaNN.tex e GabaritoNN.tex via % Preâmbulo compartilhado: MATF14, Listas de Exercícios 2026
% Usado por ListaNN.tex e GabaritoNN.tex via \input{preamble}

\usepackage[utf8]{inputenc}
\usepackage[T1]{fontenc}
\usepackage[brazil]{babel}
\usepackage{fullpage}
\usepackage{amsmath,amssymb}
\usepackage{enumitem}
\usepackage{xcolor}
\usepackage{fancyhdr}
\usepackage{titlesec}
\usepackage{listings}
\usepackage{hyperref}
\usepackage{booktabs}
\usepackage{graphicx}
\usepackage{tikz}

\definecolor{matf14azul}{HTML}{0D3B54}
\definecolor{matf14dourado}{HTML}{C9971E}

\titleformat{\section}{\normalfont\Large\bfseries\color{matf14azul}}{\thesection}{1em}{}
\titleformat{\subsection}{\normalfont\large\bfseries\color{matf14azul}}{\thesubsection}{1em}{}

\providecommand{\ListaTitulo}{Lista de Exercícios}

\setlength{\headheight}{14pt}
\pagestyle{fancy}
\fancyhf{}
\lhead{\small MATF14: Estatística Econômica I}
\rhead{\small \ListaTitulo}
\cfoot{\thepage}
\renewcommand{\headrulewidth}{0.4pt}

\lstset{
  basicstyle=\ttfamily\small,
  breaklines=true,
  frame=single,
  columns=fullflexible,
  backgroundcolor=\color{gray!8},
  commentstyle=\color{matf14dourado},
  keywordstyle=\color{matf14azul}\bfseries,
  language=R
}

\newcounter{questaocounter}
\newenvironment{questao}[1]{%
  \stepcounter{questaocounter}%
  \bigskip\noindent\textbf{Questão #1.}\ %
}{\par}

\newcommand{\cabecalho}[2]{%
  \begin{center}
    {\Large\bfseries\color{matf14azul} #1}\\[0.3em]
    {\large #2}
  \end{center}
  \vspace{0.5em}
  \noindent\rule{\linewidth}{0.6pt}
  \vspace{0.5em}
}

\newenvironment{solucao}{%
  \par\smallskip\noindent\textit{\color{matf14dourado!80!black}Solução.}\ %
}{\hfill$\blacksquare$\par}


\usepackage[utf8]{inputenc}
\usepackage[T1]{fontenc}
\usepackage[brazil]{babel}
\usepackage{fullpage}
\usepackage{amsmath,amssymb}
\usepackage{enumitem}
\usepackage{xcolor}
\usepackage{fancyhdr}
\usepackage{titlesec}
\usepackage{listings}
\usepackage{hyperref}
\usepackage{booktabs}
\usepackage{graphicx}
\usepackage{tikz}

\definecolor{matf14azul}{HTML}{0D3B54}
\definecolor{matf14dourado}{HTML}{C9971E}

\titleformat{\section}{\normalfont\Large\bfseries\color{matf14azul}}{\thesection}{1em}{}
\titleformat{\subsection}{\normalfont\large\bfseries\color{matf14azul}}{\thesubsection}{1em}{}

\providecommand{\ListaTitulo}{Lista de Exercícios}

\setlength{\headheight}{14pt}
\pagestyle{fancy}
\fancyhf{}
\lhead{\small MATF14: Estatística Econômica I}
\rhead{\small \ListaTitulo}
\cfoot{\thepage}
\renewcommand{\headrulewidth}{0.4pt}

\lstset{
  basicstyle=\ttfamily\small,
  breaklines=true,
  frame=single,
  columns=fullflexible,
  backgroundcolor=\color{gray!8},
  commentstyle=\color{matf14dourado},
  keywordstyle=\color{matf14azul}\bfseries,
  language=R
}

\newcounter{questaocounter}
\newenvironment{questao}[1]{%
  \stepcounter{questaocounter}%
  \bigskip\noindent\textbf{Questão #1.}\ %
}{\par}

\newcommand{\cabecalho}[2]{%
  \begin{center}
    {\Large\bfseries\color{matf14azul} #1}\\[0.3em]
    {\large #2}
  \end{center}
  \vspace{0.5em}
  \noindent\rule{\linewidth}{0.6pt}
  \vspace{0.5em}
}

\newenvironment{solucao}{%
  \par\smallskip\noindent\textit{\color{matf14dourado!80!black}Solução.}\ %
}{\hfill$\blacksquare$\par}


\usepackage[utf8]{inputenc}
\usepackage[T1]{fontenc}
\usepackage[brazil]{babel}
\usepackage{fullpage}
\usepackage{amsmath,amssymb}
\usepackage{enumitem}
\usepackage{xcolor}
\usepackage{fancyhdr}
\usepackage{titlesec}
\usepackage{listings}
\usepackage{hyperref}
\usepackage{booktabs}
\usepackage{graphicx}
\usepackage{tikz}

\definecolor{matf14azul}{HTML}{0D3B54}
\definecolor{matf14dourado}{HTML}{C9971E}

\titleformat{\section}{\normalfont\Large\bfseries\color{matf14azul}}{\thesection}{1em}{}
\titleformat{\subsection}{\normalfont\large\bfseries\color{matf14azul}}{\thesubsection}{1em}{}

\providecommand{\ListaTitulo}{Lista de Exercícios}

\setlength{\headheight}{14pt}
\pagestyle{fancy}
\fancyhf{}
\lhead{\small MATF14: Estatística Econômica I}
\rhead{\small \ListaTitulo}
\cfoot{\thepage}
\renewcommand{\headrulewidth}{0.4pt}

\lstset{
  basicstyle=\ttfamily\small,
  breaklines=true,
  frame=single,
  columns=fullflexible,
  backgroundcolor=\color{gray!8},
  commentstyle=\color{matf14dourado},
  keywordstyle=\color{matf14azul}\bfseries,
  language=R
}

\newcounter{questaocounter}
\newenvironment{questao}[1]{%
  \stepcounter{questaocounter}%
  \bigskip\noindent\textbf{Questão #1.}\ %
}{\par}

\newcommand{\cabecalho}[2]{%
  \begin{center}
    {\Large\bfseries\color{matf14azul} #1}\\[0.3em]
    {\large #2}
  \end{center}
  \vspace{0.5em}
  \noindent\rule{\linewidth}{0.6pt}
  \vspace{0.5em}
}

\newenvironment{solucao}{%
  \par\smallskip\noindent\textit{\color{matf14dourado!80!black}Solução.}\ %
}{\hfill$\blacksquare$\par}


\begin{document}

\cabecalho{Gabarito 09}{Variável aleatória discreta, esperança, variância e função de distribuição acumulada (Cap. 3, Aulas 17--19)}

\begin{questao}{1} \textbf{(Função de probabilidade: número de ações de um fundo)}
\begin{solucao}
\begin{enumerate}[label=(\alph*)]
  \item $0{,}5+k+0{,}2=1 \Rightarrow k=0{,}3$.
  \item $E(X) = 0(0{,}5)+1(0{,}3)+2(0{,}2) = 0+0{,}3+0{,}4=0{,}7$.
    $\mathrm{Var}(X) = (0-0{,}7)^2(0{,}5)+(1-0{,}7)^2(0{,}3)+(2-0{,}7)^2(0{,}2)$
    $=0{,}49\times0{,}5+0{,}09\times0{,}3+1{,}69\times0{,}2 = 0{,}245+0{,}027+0{,}338=0{,}61$.
  \item $E(C)=2000\times0{,}7=1400$. $\mathrm{Var}(C)=2000^2\times0{,}61=2.440.000$;
    $dp(C)=\sqrt{2.440.000}\approx 1562$.
\end{enumerate}
\end{solucao}
\end{questao}

\begin{questao}{2} \textbf{(FDA: número de atrasos de pagamento)}
\begin{solucao}
\begin{enumerate}[label=(\alph*)]
  \item $F(0)=0{,}70$; $F(1)=0{,}70+0{,}20=0{,}90$; $F(2)=0{,}90+0{,}07=0{,}97$;
    $F(3)=0{,}97+0{,}03=1{,}00$.
  \item $P(X\le1)=F(1)=0{,}90$. $P(X>1)=1-F(1)=0{,}10$. $P(X\ge2)=1-F(1)=0{,}10$ (mesmo valor de
    $P(X>1)$ pois, sendo $X$ inteira, $X\ge 2 \Leftrightarrow X>1$).
  \item Risco alto: $P(X\ge2)=0{,}10$, espera-se que 10\% da carteira seja de risco alto.
\end{enumerate}
\end{solucao}
\end{questao}

\begin{questao}{3} \textbf{(Propriedades de $E$ e $\mathrm{Var}$: taxas e impostos)}
\begin{solucao}
\begin{enumerate}[label=(\alph*)]
  \item $E(C) = 15+0{,}02\times E(X) = 15+0{,}02\times8000 = 15+160=175$.
  \item $\mathrm{Var}(C) = 0{,}02^2\times\mathrm{Var}(X) = 0{,}0004\times 3000^2 = 0{,}0004\times
    9.000.000=3600$; $dp(C)=\sqrt{3600}=60$. (Alternativamente, direto:
    $dp(C)=|0{,}02|\times dp(X)=0{,}02\times3000=60$.)
  \item Não muda o desvio padrão: a constante $b=15$ na fórmula $\mathrm{Var}(aX+b)=a^2
    \mathrm{Var}(X)$ não aparece no resultado. Remover a taxa fixa muda $E(C)$ (de 175 para 160),
    mas não muda $dp(C)$, que continua $60$, a taxa fixa desloca o custo esperado, mas não altera
    sua dispersão.
\end{enumerate}
\end{solucao}
\end{questao}

\begin{questao}{4} \textbf{(Soma de variáveis independentes: carteira de dois ativos)}
\begin{solucao}
\begin{enumerate}[label=(\alph*)]
  \item $E(W)=E(X)+E(Y)=100+150=250$. Como $X,Y$ são independentes,
    $\mathrm{Var}(W)=\mathrm{Var}(X)+\mathrm{Var}(Y)=400+900=1300$.
  \item $dp(W)=\sqrt{1300}\approx 36{,}06$. $dp(X)+dp(Y) = \sqrt{400}+\sqrt{900}=20+30=50$.
    São \textbf{diferentes} ($36{,}06 < 50$): o desvio padrão da soma é menor que a soma dos
    desvios padrão individuais, esse é precisamente o efeito de diversificação. Combinar ativos
    (mesmo sem nenhuma "mágica", só por independência) reduz o risco relativo da carteira frente à
    soma ingênua dos riscos individuais, porque as variâncias (não os desvios padrão) é que se
    somam sob independência, e a raiz quadrada da soma é sempre menor ou igual à soma das raízes.
\end{enumerate}
\end{solucao}
\end{questao}

\begin{questao}{5} \textbf{(Dedução: $\mathrm{Var}(aX+b)=a^2\mathrm{Var}(X)$)}
\begin{solucao}
\begin{enumerate}[label=(\alph*)]
  \item
  $$
  E(aX+b) = \sum_i (ax_i+b)p(x_i) = a\sum_i x_i p(x_i) + b\sum_i p(x_i) = a\mu + b\times 1 = a\mu+b,
  $$
  usando que $\sum_i x_ip(x_i)=\mu=E(X)$ (definição) e $\sum_i p(x_i)=1$ (axioma).
  \item Do item (a), $E(aX+b)=a\mu+b$. Logo,
  $$
  (aX+b) - E(aX+b) = (aX+b) - (a\mu+b) = aX - a\mu = a(X-\mu).
  $$
  \item Por definição, $\mathrm{Var}(aX+b) = E\big[((aX+b)-E(aX+b))^2\big]$. Usando o item (b),
    $(aX+b)-E(aX+b) = a(X-\mu)$, logo
  $$
  \mathrm{Var}(aX+b) = E\big[(a(X-\mu))^2\big] = E\big[a^2(X-\mu)^2\big] = a^2\,E\big[(X-\mu)^2\big]
  = a^2\,\mathrm{Var}(X),
  $$
  pois $a^2$ é uma constante que sai da esperança (soma ponderada). De fato, $b$ não aparece no
  resultado final: somar uma constante $b$ a $X$ apenas \textbf{desloca} todos os valores possíveis
  igualmente, sem alterar a distância de cada um até a nova média (que também se desloca em $b$),
  a dispersão em torno da média permanece exatamente a mesma.
\end{enumerate}
\end{solucao}
\end{questao}

\end{document}
